%% kapa: exam=1 round=32 year=2021 period=2 form=A subject=appraisal_law subject_ko=감정평가관계법규 seq=30 number=30
%% source: Q-Net 감정평가사 기출문제 (cst00309) artlSeq=5208397 file=extracted/1008/5208397/A형/2021년도 제32회 감정평가사 1차 2교시 A형.pdf
%% conversion: kordoc 4.24.51 -> tools/gen1.py (auto); review=proofread-sonnet; flags=none
\begin{question}{30}
공간정보의 구축 및 관리 등에 관한 법령상 지목변경 신청 및 축척변경에 관한 설명이다. (\quad)에 들어갈 내용으로 각각 옳은 것은?
\begin{framedtext}
○ 토지소유자는 지목변경을 할 토지가 있으면 그 사유가 발생한 날부터 (ㄱ) 이내에 지적소관청에 지목변경을 신청하여야 한다. \par
○ 지적소관청은 축척변경을 하려면 축척변경 시행지역의 토지소유자 (ㄴ) 이상의 동의를 받아야 한다. \par
\end{framedtext}
\begin{choices}
  \item ㄱ: 30일, ㄴ: 2분의 1
  \item ㄱ: 30일, ㄴ: 3분의 2
  \item ㄱ: 60일, ㄴ: 2분의 1
  \item ㄱ: 60일, ㄴ: 3분의 2
  \item ㄱ: 90일, ㄴ: 3분의 2
\end{choices}
\end{question}
\answer{4}
